\subsection{工学的意思決定プロセス}

\subsubsection{プロセス概要}

工学的意思決定は、次の流れで考えると理解しやすい。
\begin{enumerate}
  \item 本当の問題を理解する。
  \item 技術的に実現可能な候補案を洗い出す。
  \item 選択基準を定義する。
  \item 各代替案を選択基準に照らして評価する。
  \item 望ましい代替案を選ぶ。
  \item 選んだ代替案の成果を監視する。
\end{enumerate}

実務では、この手順は完全な一直線ではない。候補案を考える途中で問題理解に戻ることもあり、評価中に選択基準を見直すこともある。ただし、各ステップを飛ばすと、判断の根拠が弱くなる。

\par\smallskip
\begin{center}
\centering
\IfFileExists{assets/generated_figures/ch15/fig-ch15-03.png}{\includegraphics[width=0.96\linewidth]{assets/generated_figures/ch15/fig-ch15-03.png}}{\fbox{\parbox[c][48mm][c]{0.9\linewidth}{\centering 画像生成待ち\\工学的意思決定プロセス図}}}
\par\smallskip
\refstepcounter{figure}\label{fig:ch15-03}図 \thefigure: 工学的意思決定プロセス図
\end{center}
図 \thefigure は、工学的意思決定プロセス図を視覚的に整理したものである。
\par\smallskip

\subsubsection{本当の問題を理解する}

最良の解は、本当の問題を理解しなければ選べない。表面的な要求や解決策をそのまま受け入れると、根本原因を外した案を選ぶ危険がある。

実務では「5 つのなぜ」が有効である。なぜ必要なのかを繰り返し問うことで、要求の背後にある目的や制約を明らかにできる。また、デザイン思考の共感段階や、ビジネスモデルの確認も、問題の背景を理解する助けになる。

\subsubsection{技術的に実現可能な候補案を洗い出す}

候補案の集合に最良案が含まれていなければ、最良案を選ぶことはできない。したがって、重大な判断では創造的に候補案を探す必要がある。

技術的実現可能性を確認するには、試作やピアレビューが有効である。候補案を増やすほど最良案を含む可能性は高まるが、評価コストも増える。そのため、候補案の数は判断の重要度に応じて調整する。

\subsubsection{選択基準を定義する}

選択基準は、代替案を比べるための物差しである。金銭は重要な基準だが、唯一の基準とは限らない。たとえば環境配慮、利用者満足、セキュリティ、安全性、法令順守、従業員の負荷、ブランド信頼なども基準になり得る。

選択基準はできるだけ客観的に表す。ただし、すべてを金銭に換算できるわけではない。金銭で表せない基準は、測定不能なもの、還元不能なもの、無形のものとして扱い、複数属性意思決定で考える。

\subsubsection{各代替案を選択基準に照らして評価する}

代替案を評価するときは、同じ視点、同じ期間、同じ費用項目、同じ比較基準を使う必要がある。片方の案だけ 3 年分の費用を見て、もう片方の案だけ 5 年分の費用を見ると、不公平な比較になる。

\subsubsection{望ましい代替案を選ぶ}

金銭だけが基準であれば、現在価値、将来価値、内部収益率などで最もよい案を選ぶ。複数の基準がある場合は、複数属性意思決定技法を使う。

見積りには不確実性がある。不確実性が判断結果を変える可能性があるときは、見積り範囲を考える、感度分析を行う、最終判断を遅らせるといった対応が必要である。複数の将来結果がある場合には、リスク下の意思決定や不確実性下の意思決定を使う。

\begin{notebox}{用語の見分け方}
\term{リスク下の意思決定}{Decision Making Under Risk}は、各結果の確率をある程度割り当てられる場合である。\term{不確実性下の意思決定}{Decision Making Under Uncertainty}は、確率を割り当てられない場合である。
\end{notebox}

\subsubsection{選択した代替案の成果を監視する}

選択後は、実績を見積りと比較する。これを行わなければ、見積りが良かったのか悪かったのかが分からない。実績との差を分析することで、次回以降の見積りと意思決定を改善できる。
